\documentclass[11pt,a4paper]{article}
\usepackage[utf8]{inputenc}
\usepackage[T1]{fontenc}
\usepackage{lmodern,amsmath,amssymb,textcomp}
\usepackage[margin=24mm]{geometry}
\usepackage{longtable,booktabs,array,enumitem}
\usepackage[hidelinks]{hyperref}
\setlength{\parindent}{0pt}
\setlength{\parskip}{6pt}
\emergencystretch=3em
\hypersetup{pdfauthor={Rachel Teo, Wenyi Wang, Hamdi Abderrahmene, Alejandro Zarzuelo Urdiales}}
\begin{document}
{\LARGE\bfseries A log(5/3) uniform bound for representations as two nonnegative cubes\par}\medskip
{\small Rachel Teo\ensuremath{^1}, Wenyi Wang\ensuremath{^1}, Hamdi Abderrahmene\ensuremath{^1}, Alejandro Zarzuelo Urdiales\ensuremath{^2}\par}\medskip
{\small\textbf{Affiliations:} \href{https://sakana.ai/}{\ensuremath{^1} Sakana AI}; \href{https://rsihouse.ai/openmath}{\ensuremath{^2} Open Math}.\par}

{\small\href{https://github.com/alejandrozu/openmath-2026-judging}{Code and formal proofs}\par}

\subsection{Abstract}
For ordered nonnegative representations x\ensuremath{^3}+y\ensuremath{^3}=n, the formal endpoint gives the eventual bound exp((log(5/3)+\ensuremath{\varepsilon})log n/log log n). The author-supplied ordinary proof optimizes the same weighted ternary coding method to coefficient C=log 3\ensuremath{-}H((2\ensuremath{-}\ensuremath{\surd}2)/3)\ensuremath{-}((2\ensuremath{-}\ensuremath{\surd}2)/3)log 2\ensuremath{\approx}0.4695030882. A gcd transfer and prime cutoff remove primitive and cube-free restrictions. The stronger ordinary endpoint is not asserted as Lean-checked.

\subsection{Count and uniform theorem}
Let R(n)=\#\{(x,y)\ensuremath{\in}\ensuremath{\mathbb{N}}\ensuremath{^2}:x\ensuremath{^3}+y\ensuremath{^3}=n\}, with ordered pairs and zero coordinates included. For every \ensuremath{\varepsilon}>0 there is N such that R(n)\ensuremath{\leq}exp((log(5/3)+\ensuremath{\varepsilon})log n/log log n) for every n\ensuremath{\geq}N. No cubefree-target restriction or gcd-one restriction appears in the final theorem.

The coefficient log(5/3)\ensuremath{\approx}0.510826 improves the earlier same-family log(7/4)\ensuremath{\approx}0.559616 endpoint. The located approximately0.556477 comparison applies only to cubefree targets. It therefore cannot be substituted as a stronger all-target theorem. The result still falls far short of a fixed polylogarithmic bound, which remains the original target.

\subsection{Primitive ratios as a separated ternary code}
For a primitive positive solution, the ratio of its coordinates is a unit modulo the relevant prime powers, with cube equal to minus one. The unit cube-root lemma gives at most three choices at each prime-power coordinate, including the ramified prime three. Distinct solutions are encoded by distinct ternary words.

If two encoded words agree at selected prime-power coordinates, those moduli divide a cross determinant of the two solutions. An upper bound for the determinant and the target-size product force a weighted separation property. This converts an arithmetic representation problem into a weighted code problem without assuming the desired representation bound.

Elias/Plotkin localization bounds the size of a separated code in a weighted ball. An exact binomial block certificate uses blocks of 500 coordinates and support 90, producing the rational base 5/3 with an explicit finite multiplicative constant. The constant is harmless for the eventual leading exponent, but its proof is retained rather than replaced by a decimal numerical guess.

\subsection{From primitive pairs to the canonical count}
A general positive pair has gcd d and a primitive pair representing n/d\ensuremath{^3}. The source proves the exact gcd-fibre identity and sums over positive cube divisors d. Coordinatewise multiplicative weights bound this sum uniformly across targets, including those with repeated prime factors.

The prime cutoff and elementary log budgets control the resulting product. They reuse background modules from the k\ensuremath{\sigma}(k) development, while the cubic encoding and local-factor estimates are specific to this result. Finally at most two boundary pairs have a zero coordinate, and their contribution is absorbed. A semantic bridge identifies the resulting count with the original ordered cube-value representation definition; it does not import the open conjecture as a hypothesis.

\subsection{Formal scope and prior work}
The checked source snapshot contains reconstructed arithmetic variants, identified in the source ledger. It verifies those delivered variants; it does not claim recovery of the earlier lost bytes. The selected final endpoints use standard kernel axioms.

Merikoski's ratio-root and determinant methods and classical coding/entropy estimates are prior ingredients. The refinements are the exact 5/3 weighted-code constant and the separately justified ordinary coefficient 0.4695030882.... Their all-target transfer is a quantitative partial result; the polylogarithmic target and a signed-integer representation theorem are not proved here.

The weighted-distance condition, finite block inequality and gcd-fibre identity give the proof chain from primitive representations to the canonical count. The coefficient is justified by exact finite inequalities, not a floating-point estimate.

\subsection{Formal verification}
Lean 4.33.1 checked all 31 delivered source modules and both selected declarations canonicalCubeCount\_eventually\_exp and canonicalCubeCount\_eventual\_threshold. They use only standard kernel axioms and prove the all-target log(5/3) bound. The sharper ordinary coefficient .46950308823856934 has a separate proof and is not asserted as a selected Lean endpoint. Pinned source identities and audit records are supplied in the companion repository.

\begin{samepage}
\subsection{Erdos829FiveThirds.canonicalCubeCount\_eventually\_exp}
\begin{quote}\small\ttfamily theorem canonicalCubeCount\_eventually\_exp \{\ensuremath{\varepsilon} : \ensuremath{\mathbb{R}}\} (h\ensuremath{\varepsilon} : 0<\ensuremath{\varepsilon}) :\newline     \ensuremath{\forall}\ensuremath{^{f}} n : \ensuremath{\mathbb{N}} in atTop, (canonicalCubeCount n : \ensuremath{\mathbb{R}}) \ensuremath{\leq}\newline       Real.exp ((Real.log (5/3 : \ensuremath{\mathbb{R}})+\ensuremath{\varepsilon})*Real.log (n : \ensuremath{\mathbb{R}})/Real.log (Real.log (n : \ensuremath{\mathbb{R}})))\end{quote}
{\small Source: proofs/erdos829-log-five-thirds/FiveThirdsCanonical.lean:8.\par}

{\small Source SHA-256: 51541747f8104588463d8b7b051d69692876cee605c68a2f50df9da132e6a006\par}

\end{samepage}
\begin{samepage}
\subsection{Erdos829FiveThirds.canonicalCubeCount\_eventual\_threshold}
\begin{quote}\small\ttfamily theorem canonicalCubeCount\_eventual\_threshold \{\ensuremath{\varepsilon} : \ensuremath{\mathbb{R}}\} (h\ensuremath{\varepsilon} : 0<\ensuremath{\varepsilon}) :\newline     \ensuremath{\exists} N : \ensuremath{\mathbb{N}}, \ensuremath{\forall} n : \ensuremath{\mathbb{N}}, N\ensuremath{\leq}n \ensuremath{\to} (canonicalCubeCount n : \ensuremath{\mathbb{R}}) \ensuremath{\leq}\newline       Real.exp ((Real.log (5/3 : \ensuremath{\mathbb{R}})+\ensuremath{\varepsilon})*Real.log (n : \ensuremath{\mathbb{R}})/Real.log (Real.log (n : \ensuremath{\mathbb{R}})))\end{quote}
{\small Source: proofs/erdos829-log-five-thirds/FiveThirdsCanonical.lean:14.\par}

{\small Source SHA-256: 51541747f8104588463d8b7b051d69692876cee605c68a2f50df9da132e6a006\par}

{\small\textbf{Sources and attribution:} \href{https://www.erdosproblems.com/829}{1. www.erdosproblems.com / 829}; \href{https://github.com/google-deepmind/formal-conjectures/blob/main/FormalConjectures/ErdosProblems/829.lean}{2. google-deepmind/formal-conjectures / 829.lean}; \href{https://github.com/google-deepmind/formal-conjectures/blob/main/FormalConjecturesForMathlib/Combinatorics/Additive/Convolution.lean}{3. google-deepmind/formal-conjectures / Convolution.lean}; \href{https://arxiv.org/html/2112.03617}{4. arXiv source: 2112.03617}; \href{https://arxiv.org/html/2207.12487v2}{5. arXiv source: 2207.12487v2}; \href{https://github.com/alejandrozu/ulam-opdp-difficulty-atlas}{6. alejandrozu/ulam-opdp-difficulty-atlas}; \href{https://github.com/alejandrozu/openmath-2026-judging/blob/d0ed487f487fa7ec814ff705ab55249983fe8707/OpenMath-Judging/review/entries/E02-erdos829-log-five-thirds.txt}{7. alejandrozu/openmath-2026-judging / E02-erdos829-log-five-thirds.txt}; \href{https://github.com/alejandrozu/openmath-2026-judging}{Proof source repository}.\par}

\end{samepage}
{\small \textbf{AI assistance.} Codex assisted literature checking, editorial preparation and analysis of the supplied formal artifacts. The human authors are responsible for checking the final mathematical claims, references and attribution.\par}

\end{document}
